\documentclass[11pt,a4paper]{article}
\usepackage[margin=24mm]{geometry}
\usepackage{amsmath,amssymb,amsthm,booktabs,longtable,array,hyperref}
\hypersetup{hidelinks}
\emergencystretch=3em
\newcommand{\C}{\textbf{(c)}\ }
\newcommand{\D}{\textbf{(d)}\ }
\newcommand{\scr}[1]{\texttt{uniform\_\allowbreak limits\_\allowbreak #1.py}}
\newcommand{\file}[1]{\nolinkurl{#1}}
\newcolumntype{L}[1]{>{\raggedright\arraybackslash}p{#1}}
\newtheorem{proposition}{Conditional proposition}
\title{Uniform Toda limits and robust subleading asymptotics:\\boundary layers, endpoint matching, and remaining gaps}
\author{Research note for Gustav Delius}
\date{4 November 2026}
\begin{document}
\maketitle
\begin{abstract}
\C The bulk kernel expansion misses a delta-like contribution of mass
$\pi h$. On fixed interior sets and for uniformly Lipschitz test functions,
its inclusion leaves an $O(h^2\log(1/h))$ operator remainder. This is not a
uniform inverse estimate for the singular Bethe density.
\C In the corrected half-line model, an explicit matched-endpoint argument,
conditional on stated differentiated asymptotic bounds, protects the constant
and $Y^{-1}$ two-body coefficients. Independent endpoint-resolved numerical
convolution tests their negative subleading signs without prescribing a leading
fit intercept. The classical Taylor term $(\log Y)^2/Y$ is real; its absence
from the complete half-line phase at this order is a cancellation, not permission
to omit it from individual pieces.
\C Read-only refits reproduce all three failures of the Phase 2d step bound
and the Phase 2f sign reversal. The one-body nested-limit coefficient remains
compatible with $1/(2\pi)$ in the published deterministic envelope.
\D Uniform finite-sea-to-half-line resolvent control, interchange with the
fixed-nonzero-momentum limit, and the physical two-body quantum frame are
not established. This report makes partial progress on open problem 3; it does
not declare it solved.
\end{abstract}

\section{Scope, precedence and conventions}
Labels mean (a) proved in literature, (b) literature heuristic, (c) project
calculation or explicitly conditional implication, and (d) unproved hypothesis
or identification. Numerical (c) statements concern the sampled problem only.
The controlling sources are the README status and master-report conclusions,
\file{errata.tex}, \file{audit_phase2i.tex}, \file{audit_frame_dictionary.tex},
and the one-body precision note, with the later scheme audit taking precedence
on sound normalisation. The Phase 2h independent checker and Phase 2d/2f
implementations and caches were also read. No external paper, abstract,
DOI or arXiv entry was newly read or cited. Internal references below identify
actual files, not independent literature proofs.

\C The calculation retains the existing Toda extraction: $a=1$,
$b=e^{-1/2}$, $h=\hbar_T=3b\hbar_B$ and
\begin{equation}
 P_0(\kappa)=4b(\kappa\cosh\kappa-\sinh\kappa),\quad
 E_0(\kappa)=2b^2(\sinh\kappa\cosh\kappa-\kappa),\quad
 Y=b\kappa^2/h.
 \label{eq:scales}
\end{equation}
The inherited continuum dictionary is the \emph{bad} Boussinesq scaling,
$\kappa=\sqrt3\beta$ to leading order, $k=\sinh\beta$, $v=\epsilon\cosh\beta$,
$u_{tt}=u_{xx}+u_{xxxx}+3(u^2)_{xx}$, rather than a newly chosen good-equation
regularisation. Its classical charge conversion is
$P_B=\overline P/(6\sqrt3 b)$ and $E_B=\overline E/(6\sqrt3 b^2)$.
No change to the existing sound line, quantum frame, or regulator is made.
Checks: \scr{analytic}, \scr{scan}, \scr{refit}.

\C At every finite $h$, a prescribed momentum is enforced by solving
$\overline P(p)=P_0(\kappa)$ before evaluating energy or phase.
``Fixed dressed momentum'' in an edge scan means this inversion at each
point; the target itself shrinks as $h^{3/2}$ when $Y$ is fixed.
It does not mean a fixed nonzero momentum while $h$ varies.
Indeed $P_0'=4b\kappa\sinh\kappa>0$ and
$P_0\sim4b\kappa^3/3$ (\scr{analytic}). At fixed units, strictly nonzero fixed
$P_0$ fixes $\kappa$; a subsequent small-amplitude limit is a separate operation.

\section{What the kernel does and does not control}
\C Define the finite-sea problem exactly as in the audited precision package:
\begin{align}
 K_h(d)&=2\Re\psi(1+id/h)+2\log h,\qquad
 \Theta_h(d)=2h\Im\log\Gamma(1+id/h)+2d\log h,\\
 \mathcal A_h f&=2\pi h f-\int_{-q}^qK_h(x-z)f(z)\,dz,
 \quad \mathcal A_h\rho=1,\quad \mathcal A_h\eta=x,
 \quad\int_{-q}^q\rho=1.\label{eq:sea}
\end{align}
The log-Gamma phase is its continuous branch on the real argument line.
These definitions and the endpoint counting check are implemented in
\scr{scan}; they are not modified to improve agreement.

\paragraph{A uniform direct-operator estimate (c; \scr{analytic}).}
Put
\begin{equation}
 R(x)=2\Re\psi(1+ix)-2\log|x|,
 \qquad K_h(d)=2\log|d|+R(d/h).
\end{equation}
The primitive on $x>0$ and its tail give
\begin{align}
 \int_0^L R(x)\,dx
 &=2\Im\log\Gamma(1+iL)-2L\log L+2L,\\
 R(x)&=\frac1{6x^2}+\frac1{60x^4}+O(x^{-6}),\qquad
 \int_{\mathbb R}R(x)\,dx=\pi.\label{eq:mass}
\end{align}
Near zero $R(x)=-2\log|x|-2\gamma+O(x^2)$, so it is absolutely integrable.
Thus, for $x$ a fixed distance $\delta>0$ from the endpoints, and $f$ with
uniformly bounded Lipschitz norm,
\begin{align}
 \int_{-q}^q R((x-z)/h)f(z)\,dz
 &=\pi h f(x)+O_{q,\delta}(h^2\log(1/h)\|f\|_{\rm Lip}),\\
 \mathcal A_h f
 &=-\int_{-q}^q2\log|x-z|f(z)\,dz+\pi h f(x)
   +O_{q,\delta}(h^2\log(1/h)\|f\|_{\rm Lip}).\label{eq:operator}
\end{align}
To prove the estimate, subtract $f(x)$ inside the integral, use
$|f(x-hs)-f(x)|\leq h|s|\operatorname{Lip}(f)$ and
$\int^{O(1/h)} |sR(s)|\,ds=O(\log(1/h))$; the missing tails multiplying
$f(x)$ are $O(h^2/\delta)$. This uses absolute-value estimates and no assumed
positivity of $R$. The logarithm is not merely formal: the test function
$|z|$ at $x=0$ gives
\begin{equation}
 \int_{-1}^1R(z/h)|z|\,dz
 =\frac{h^2}{3}\log(1/h)+O(h^2).\label{eq:cusp}
\end{equation}
Direct 70-digit integration verifies the primitive, mass tail and logarithmic
moment increment. The asymptotic equalities follow from the displayed
primitive and digamma expansion, not from extrapolating those numerical tests.

\C At fixed $d\ne0$, $R(d/h)=O(h^2)$, whereas on $d=hu$ it is $R(u)$.
The pointwise expansion therefore cannot be inserted termwise across the sea
without matching the diagonal region. In particular $\mathcal A_h$ has a
$\pi h$ local term, not the $2\pi h$ suggested by discarding that region.
\D Equation~\eqref{eq:operator} is not an estimate for
$\mathcal A_h^{-1}$ on the endpoint-singular densities or their $q$ derivatives.
No bound on that inverse, the moving endpoint, or its differentiated error is
proved here. Consequently $h\log h$ corrections in observables divided by $h$
are admissible fit alternatives; this calculation neither forces nor excludes
them for the actual dressed phase. Smooth interior data can have additional
cancellations not available for a generic Lipschitz test function.

\section{Corrected half-line asymptotics and one-body limits}
\paragraph{Precisely which half-line problem (c, conditional; \scr{analytic}).}
Use the corrected factor and normalization:
\begin{equation}
 g(s)=s^{-1/2}\Gamma(1+s/(2\pi))
 \exp\{-s\log[s/(2\pi e)]/(2\pi)\},\qquad
 \mathcal L F=\sqrt\pi\,g(s)/s^2,\quad F(u)=\tfrac43 y(u)^{3/2}.
 \label{eq:wh}
\end{equation}
The branch is principal for $\Re s>0$; define $B=F'$ and $f=F''$ on $u>0$.
Then $\mathcal L B=\sqrt\pi g/s$ and
$\mathcal L f=\sqrt\pi(g-1)$, since $B(0^+)=\sqrt\pi$.
The nonzero-frequency product is
$2\pi g(i\omega)g(-i\omega)=2\pi/(1-e^{-|\omega|})$.
This is a corrected half-line model, not a proof that the zero-frequency
finite-sea matching has the required uniform remainder.

\C Formal inverse Laplace expansion and reversion, checked symbolically
through cubic order in $s$ and at 70 digits, give
\begin{align}
 u(y)&=y-A\log y+c+\frac d y+\frac e{y^2}+\cdots,\label{eq:u}\\
 A&=\frac1{2\pi},\quad c=\frac{1-\log(8\pi)}{2\pi},\quad
 d=\frac1{48}-\frac1{4\pi^2},\quad
 e=\frac{4+6\zeta(3)-\pi^2}{96\pi^3}.\nonumber
\end{align}
There is no $(\log y)/y$, $(\log y)^2/y^2$ or $(\log y)/y^2$ in this
\emph{inverse} series. This does not remove logarithms from $y(u)$.
The coefficients are $d=-0.00449696257725111$ and
$e=0.000451097178817793$. Their interpretation as coefficients of the
finite-sea limit remains conditional on matching.

\C With $J(u)=\int_0^uF(w)\,dw$ and the inherited harmonic sound shift,
\begin{align}
 C_E(y)&=\frac3{4\pi}+\frac{3J(u(y))}{4F(u(y))}-\frac{3y}{10},
 &C_v(y)&=\frac3{4\pi}+\frac y2(u'(y)-1),\label{eq:CE}\\
 C_E(y)&=\frac1{2\pi}-\frac{3d}{2y}
       +\frac{9K}{16y^{3/2}}+\frac{3e}{y^2}+\cdots,
 &C_v(y)&=\frac1{2\pi}-\frac d{2y}-\frac e{y^2}+\cdots.\label{eq:tails}
\end{align}
Here $K$ is the possible integration constant when the large-$y$ derivative
series is integrated in isolation. Keeping it in finite-data fits avoids
assuming an unproved endpoint sum rule. The sound coefficient follows from
$e_{\rm vac}(a,h)=b^2+2hb/\pi+O(h^2)$ and
$v_s=a\sqrt{\partial_a^2e_{\rm vac}}=ab+ah/(4\pi)+O(h^2)$.
Checks: \scr{analytic}, \scr{matched}, \scr{refit}.

\paragraph{A concrete order-of-limits distinction (c, conditional on this model).}
For $y\downarrow0$, $F\sim\sqrt\pi u$, so $u\sim4y^{3/2}/(3\sqrt\pi)$.
Consequently both quantities in \eqref{eq:CE} tend to $3/(4\pi)$, whereas
for $y\to\infty$ they tend to $1/(2\pi)$. At finite $h$, taking $P\to0$
first in $3[\overline E(P)-E_0(P)]/(hP)$ gives $3(v_s-b)/h$.
Thus the phonon-first limit is $3/(4\pi)$, differing from the nested
soliton-edge coefficient by $1/(4\pi)$. These two regimes must not be
called a single path-independent joint limit. Checks: \scr{analytic}.
\D This distinction does not evaluate $\lim_{\kappa\to0}\lim_{h\to0}$
at fixed nonzero $\kappa$ in the inner limit; that is precisely the missing
interchange addressed numerically below.

\section{Matched endpoint analysis of the two-body half-line phase}
\C The two-term phase tested here, with its inherited orientation and
normalization, is
\begin{equation}
 Q(u,v)=2\int_0^v f(u-v+z)B(z)\,dz,\quad u>v>0,
 \qquad R(Y,r)=Q(u(Y),u(r^2Y))-\Phi_{\rm cl}(Y,r),\label{eq:Q}
\end{equation}
where
$\Phi_{\rm cl}(Y,r)=2Y(1-r^2)\log[(1-r)/(1+r)]+4rY$.
This is the corrected convolution tested in the Phase 2i audit.
It is not substituted for an antisymmetric physical scattering phase.
The finite-sea counterpart is exactly the historical two-term expression,
with its classical frame subtraction and full classical lattice phase removed.
At leading edge order that frame subtraction removes $4rY$; the finite-sea
classical reference consequently tends to $2Y(1-r^2)\log[(1-r)/(1+r)]$.
The resulting residual is the same $R$ as in \eqref{eq:Q}; the two
unsubtracted phases are not being equated.
Checks: \scr{scan}, \scr{wh}, \scr{matched}, \scr{refit}.

\begin{proposition}[matched half-line implication, (c)]
Assume $u$ is increasing and regular on the open positive axis, obeys the
small-$y$ behavior stated above, and for $y\geq1$ the remainder after
$y-A\log y+c+d/y$ and its first three derivatives are bounded by constants
times $y^{-2-j}\log^m(2+y)$, $j=0,1,2,3$. For $r$ in a fixed compact subset
of $(0,1)$, the convolution \eqref{eq:Q} then has
\begin{align}
 R(Y,r)&=c_0(r)+\frac{c_1(r)}Y
       +O\bigl(Y^{-3/2}\log^M Y\bigr),\label{eq:result}\\
 c_0(r)&=-\frac2\pi[\operatorname{Li}_2(r)-\operatorname{Li}_2(-r)-2r],\nonumber\\
 c_1(r)&=-\frac{r^3}{6(1-r^2)}
 -2d\left[\frac{1-r^2}{r^2}\log\frac{1-r}{1+r}+\frac2r\right].\label{eq:c1}
\end{align}
Here $M$ is a finite logarithmic power depending on the remainder hypotheses;
no uniformity as $r\to0$ or $r\to1$ is asserted.
\end{proposition}
\C Checks for the algebra, exact transformed measure and coefficient
integrals are in \scr{matched}; the implication's endpoint estimates follow
below. \D The differentiated bounds are hypotheses: symbolic reversion and
numerical inversion are evidence for them in \eqref{eq:wh}, not a proved
finite-sea resolvent theorem.

\paragraph{Derivation and the previously missing endpoint estimate (c).}
Changing variables $z=u(x)$ makes $B(z)\,dz=2\sqrt{x}\,dx$ exact. Put
$x=Yt^2$, $\sigma=1-r^2+t^2$ and
$\Delta_0=-A\log[t^2/(r^2\sigma)]$. Differentiating $F(u(x))$ twice yields
\begin{equation}
 f(u(x))=\frac1{\sqrt{x}\,u'(x)^2}
 -\frac{2\sqrt{x}\,u''(x)}{u'(x)^3}
 =x^{-1/2}\left[1-\frac{2d+3A^2}{x^2}+\cdots\right].
\end{equation}
The $x^{-1}$ relative term cancels. At every fixed $t>0$, solving for the
argument's dressed coordinate gives the balanced displacement
\begin{equation}
 \delta(Y)-\delta(r^2Y)+\delta(Yt^2)-\delta(Y\sigma),
 \qquad\delta(y)=u(y)-y.
\end{equation}
Its leading term is $\Delta_0$: both $\log Y$ and $c$ cancel.
The constant and $Y^{-1}$ integrands, after the classical term is removed, are
\begin{align}
 I_0(t)&=-4t^2\Delta_0/\sigma^{3/2},\\
 I_1(t)&=\frac8{\sigma^{5/2}}
 \left[t^2\left(-2d-3A^2-\frac{A\Delta_0}{2}
                       +\frac{3\Delta_0^2}{8}\right)
 -\frac{d\sigma}{2}\{1+t^2(1-r^{-2}-\sigma^{-1})\}\right].
\end{align}
Direct independent 60-digit quadrature of $I_0,I_1$ agrees with
\eqref{eq:c1} and $c_0$ at five ratios from $0.2$ to $0.8$, to below $10^{-54}$.

\C Split the exact integral at $t=L/\sqrt Y$, for a fixed sufficiently large
$L$. On the inner interval the exact displacement and its truncated logarithm
are bounded by a power of $\log Y$ plus the integrable $d/(Yt^2)$ term.
Their contribution to the residual is bounded by a constant times
$Y^{-3/2}\log^M Y$: $\int_0^{L/\sqrt Y}t^2dt=O(Y^{-3/2})$ and the last
term contributes $Y^{-1}\int_0^{L/\sqrt Y}dt=O(Y^{-3/2})$.
On the outer interval the worst neglected displacement contributes
$Y^{-2}\int_{L/\sqrt Y}^r t^{-2}\log^m(2+Yt^2)dt
=O(Y^{-3/2}\log^M Y)$; other Taylor remainders are no larger.
The argument of $f$ stays of order $Y$ because $1-r^2$ is bounded below.
This supplies an integrable endpoint bound at the first two orders under the
stated hypotheses. It does not justify the stronger historical
$O(Y^{-2})$ remainder without further moment matching.

\paragraph{The logarithmic Taylor term is not zero (c; \scr{analytic}).}
For displacements $d_1,d_2$ of $(Y,r^2Y)$, the quadratic Taylor term of
$\Phi_{\rm cl}$ is
\begin{equation}
 \frac{(d_2-r^2d_1)^2}{r(1-r^2)Y}.
\end{equation}
Taking $d_1=-A\log Y+c$, $d_2=d_1-\log r/\pi$ gives the nonzero coefficient
$(1-r^2)/(4\pi^2r)$ of $(\log Y)^2/Y$.
It cancels only after the other half-line contributions and exact dressed
measure are included. The proposition explains that cancellation at order
$Y^{-1}$ conditionally, without discarding the small endpoint.

\paragraph{A further endpoint moment, not a finite-sea theorem (c).}
Define the convergent moment
\begin{equation}
 \mu=\int_0^\infty\sqrt{x}
 [u(x)-x+A\log x-c-d/x],dx.\label{eq:mu}
\end{equation}
Integration by parts gives $K=-2\mu$ in \eqref{eq:tails}.
The first fractional endpoint contribution in the phase is correspondingly
$K/(1-r^2)^{3/2}\,Y^{-3/2}$.
For the formal transform \eqref{eq:wh}, $\mathcal LJ=\sqrt\pi g/s^3$
has only half-integer algebraic powers (and their logarithms); its formal
large-$y$ series contains no constant $K$. This predicts $\mu=K=0$.
\scr{matched} checks the finite-cutoff moment independently against direct
quadrature at $y=4$, and tests its large-cutoff tail
$\mu_{[0,Y]}=-2e/\sqrt Y+O(Y^{-3/2}\log^mY)$ through $Y=4096$.
\D A full contour remainder proof, or the extra matching needed to promote
an $O(Y^{-2})$ remainder to the finite-sea theory, is not supplied.
Fractional and logarithmic alternatives therefore remain in the fit ensemble.

\section{Numerical results and separated uncertainties}
\paragraph{One-body: a renewed fit to the saved finite-sea data (c).}
The constant remains free. Six scaled $h$ models use polynomial degrees
2--5, the last 4--8 points, and both $h\log h$ and $h^2\log h$ alternatives.
After subtracting only the derived $-3d/(2y)$ tail, sixteen amplitude fits
vary windows within $4\leq y\leq64$ and allow powers, $y^{-3/2}$, or
$y^{-2}\log y$. The central fit and its separately propagated diagnostics are

\[C_T=0.15915494136\ \pm\ 8.23\times10^{-6}.\]
\begin{center}\begin{tabular}{lr}\toprule Component & Size\\\midrule
Mesh/roundoff & $5.68\times10^{-8}$\\
$h$-basis/window spread & $7.70\times10^{-6}$\\
Amplitude-basis/window spread & $4.73\times10^{-7}$\\
\bottomrule\end{tabular}\end{center}
This reproduces the precision note's envelope with a new read-only fit.
It is not a new finite-sea dataset or a confidence interval. Absolute
regression weights propagate mesh differences plus a $2\times10^{-10}$ floor;
model spreads are added linearly. No allowance bounds failure of the edge
ansatz or the physical charge dictionary. Check: \scr{refit}.

\paragraph{Phase 2d: the bound is false, not a failed convergence test (c).}
The saved raw phase differences reproduce these violations:
\begin{center}\begin{tabular}{rrrr}\toprule
$(y_1,y_2)$ & $h$ before & $h$ after & Absolute change\\\midrule
$(3,1)$ & 0.0100 & 0.0050 & 0.006326022\\
$(6,1)$ & 0.0100 & 0.0050 & 0.009016413\\
$(6,1)$ & 0.0050 & 0.0025 & 0.004511945\\
\bottomrule\end{tabular}\end{center}
For $(6,1)$ the ratio of the two successive changes is $0.500415$,
compatible with a leading linear-$h$ correction over these three points.
Three points neither establish that asymptotic law nor exclude $h\log h$.
The original $0.004$ assertion remains false; no replacement universal bound
is asserted. Check: \scr{refit}.

\paragraph{Phase 2f: robust saved-data refits (c).}
First subtract the full classical lattice phase, evaluated at 60 digits using
the hyperbolic cross ratio, instead of only its edge limit:
\begin{equation}
 \Phi_{\rm lat}/h=\frac{2b(\cosh\kappa_1-\cosh\kappa_2)\log A_T
       +4b(\kappa_2\sinh\kappa_1-\kappa_1\sinh\kappa_2)}h,
 \quad A_T=\frac{\sinh^2[(\kappa_1-\kappa_2)/2]}
                  {\sinh^2[(\kappa_1+\kappa_2)/2]}.
\end{equation}
All smallest-3, smallest-4 and all-5 point sets are tested with
$(1,h,h^2)$, $(1,h,h^2,h^3)$, $(1,h,h\log h)$,
$(1,h,h^2,h^2\log h)$ and $(1,h,h^2,h\log h)$ whenever determined.
For the subsequent $Y$ fit, both $\{8,16,32\}$ and $\{4,8,16,32\}$ are used,
with a free constant and $1/Y$ coefficient, plus respectively
$1/Y^2$, $\log Y/Y$ or $(\log Y)^2/Y$. These are sensitivity fits;
exact interpolation of a sparse dataset does not select a basis.
\begin{center}\begin{tabular}{rrr}\toprule
$r$ & Range of free $c_0$ & Range of free $1/Y$ coefficient\\\midrule
0.4 & $[-0.010315,-0.005614]$ & $[-0.050095,0.028953]$\\
0.5 & $[-0.020902,-0.012102]$ & $[-0.100731,0.042933]$\\
0.6 & $[-0.037707,-0.023335]$ & $[-0.178826,0.054938]$\\
\bottomrule\end{tabular}\end{center}
The $1/Y$ coefficient here is basis-dependent when logarithms are allowed;
these ranges must not be read as estimates of a unique physical coefficient.
They expose what the saved data cannot resolve. The original, unsubtracted
Phase 2f two-point slopes also reverse sign between the recorded quadratic
and $h\log h$ extrapolations:
\begin{center}\begin{tabular}{rrrr}\toprule
$r$ & Quadratic & Cubic & $h\log h$\\\midrule
0.4 & -0.0070801 & -0.0072612 & 0.0091481\\
0.5 & -0.0209913 & -0.0213631 & 0.0178186\\
0.6 & -0.0482194 & -0.0489074 & 0.0303174\\
\bottomrule\end{tabular}\end{center}
The saved mesh difference at $h=0.01$ is $2.13\times10^{-12}$;
that one refinement cannot bound mesh errors at every other $h$.
Checks and complete fit coefficients/weights: \scr{refit} and its JSON.

\paragraph{Endpoint-resolved half-line phase (c, conditional model).}
Fresh convolution data use $Y=16,32,64,128,256,512$ and
$r=0.4,0.5,0.6$. Neither $c_0$ nor $c_1$ is fixed in the fits.
The central basis is $(1,Y^{-1},Y^{-2},Y^{-3})$ on $Y\geq32$.
Alternatives start at $16$ and include $Y^{-3/2}$, or both
$(\log Y)/Y$ and $(\log Y)^2/Y$ with a $Y^{-2}$ tail.
\begin{center}\small\begin{tabular}{rrrrr}\toprule
$r$ & Free $c_1$ & Model spread & Quadrature diagnostic & Formal $c_1$\\\midrule
0.4 & -0.0077363875 & $5.36\times10^{-4}$ & $1.10\times10^{-8}$ & -0.0077365880\\
0.5 & -0.0214445505 & $3.36\times10^{-5}$ & $3.97\times10^{-8}$ & -0.0214445873\\
0.6 & -0.0484360177 & $5.36\times10^{-4}$ & $5.32\times10^{-8}$ & -0.0484359877\\
\bottomrule\end{tabular}\end{center}
The largest raw quadrature difference is $7.23\times10^{-13}$; the independent inversion-method difference at the sampled arguments is below $2.70\times10^{-41}$.
Absolute fit weights propagate the quadrature differences plus a $10^{-28}$
arithmetic floor. The much larger model spread is the honest limitation
of the subleading extraction. Every tested coefficient is negative, and
\eqref{eq:c1} lies inside these diagnostic envelopes. Tiny fitted logarithmic
coefficients are not a numerical proof of exact zeros. Checks: \scr{wh},
\scr{refit}. A conditional half-line sign does not resolve the sign from
finite-sea data alone.

\paragraph{Fresh fixed-$P$ and growing-$Y$ finite-sea checks (c).}
The base meshes at $h=0.02,0.01,0.005,0.0025$ have
$N=700,1200,2200,4200$ nodes. Separate refinements at the last three values
use $N=1500,2750,5250$. Ground-state normalization, edge counting,
fixed-momentum inversion, positive dressed derivative and the exact
thermodynamic derivative identity are asserted. The same $h$-basis ensemble
as above is used on the four available points. The central fixed-$\kappa$
estimate is the quadratic fit through the finest three values.
The phase column is the classical-subtracted two-term phase at
$\kappa_2/\kappa_1=0.5$.
\begin{center}\small\begin{tabular}{rrrrrr}\toprule
$\kappa_1$ & $C(\kappa_1)$ & $h$ spread & Phase & $h$ spread & Max. mesh diagnostic\\\midrule
0.12 & 0.15915758 & $1.41\times10^{-3}$ & -0.0194700 & $2.28\times10^{-2}$ & $4.52\times10^{-9}$\\
0.18 & 0.15878874 & $4.84\times10^{-4}$ & -0.0212320 & $3.48\times10^{-3}$ & $3.53\times10^{-9}$\\
0.25 & 0.15842942 & $8.29\times10^{-5}$ & -0.0217711 & $8.90\times10^{-4}$ & $3.83\times10^{-9}$\\
0.35 & 0.15773618 & $1.67\times10^{-5}$ & -0.0238967 & $1.71\times10^{-5}$ & $4.99\times10^{-9}$\\
0.50 & 0.15628198 & $3.32\times10^{-6}$ & -0.0282305 & $2.17\times10^{-6}$ & $7.57\times10^{-9}$\\
0.70 & 0.15360594 & $6.36\times10^{-7}$ & -0.0357984 & $4.63\times10^{-7}$ & $1.12\times10^{-8}$\\
\bottomrule\end{tabular}\end{center}
The largest available mesh difference plus a $2\times10^{-10}$ floor is
propagated with the largest absolute fit-weight norm. This is an empirical
mesh diagnostic, including an assumption that the unrefined coarsest point
is no worse; it is not a rigorous discretization bound. The large spreads
at the smallest amplitudes show that these data do not yet isolate the
fixed-$P$ semiclassical limit there. No small-$\kappa$ extrapolation from
this table is advertised as a uniform result.

For $r=0.5$, the fresh fixed-$Y$ phase comparison is:
\begin{center}\begin{tabular}{rrrrr}\toprule
$Y$ & Finite-sea intercept & $h$ spread & Mesh diagnostic & WH discrepancy\\\midrule
16 & -0.020854748 & $2.85\times10^{-4}$ & $3.88\times10^{-9}$ & $1.50\times10^{-7}$\\
32 & -0.020187307 & $9.79\times10^{-4}$ & $4.65\times10^{-9}$ & $2.06\times10^{-6}$\\
64 & -0.019876722 & $2.94\times10^{-3}$ & $6.21\times10^{-9}$ & $2.63\times10^{-5}$\\
128 & -0.019970355 & $6.81\times10^{-3}$ & $9.28\times10^{-9}$ & $2.87\times10^{-4}$\\
\bottomrule\end{tabular}\end{center}
The central intercept is cubic in $h$. The growth of the extrapolation
sensitivity with $Y$ matters even when mesh differences are tiny. It is
not a controlled expansion along $Y\propto h^{-1}$. Checks:
\scr{scan}, \scr{refit}.


\paragraph{Reproduction and arithmetic scope.}
\C All five scripts are new and write only \file{uniform_limits_*.json}.
\scr{analytic} uses SymPy and 70-digit mpmath; \scr{matched} and
\scr{refit} use 60 digits. \scr{wh} uses 40 digits, an independent inverse
Laplace method at sampled arguments, and 12-/20-node Gauss quadrature on
explicit geometric panels beginning at zero. That panel construction keeps
the $z=O(1)$ region as $Y$ grows. The dense finite-sea solver in
\scr{scan} uses double precision, with 60-digit classical energy subtraction;
higher-precision fitting does not upgrade those data.
The solver algorithm is adapted from the audited precision implementation,
not claimed as a second independent discretisation. Existing modules are not
imported or executed. Input hashes in the refit cache guard against concurrent
changes during that run.

\C To reproduce, run only these files from \file{phase2/scripts/}:
\begin{verbatim}
python3 uniform_limits_analytic.py
python3 uniform_limits_scan.py
python3 uniform_limits_wh.py
python3 uniform_limits_matched.py
python3 uniform_limits_refit.py
\end{verbatim}
The scan, WH and matched-moment scripts validate their completed numerical caches by default;
\file{--refresh} regenerates them. The two main numerical timings are recorded in their
caches. The report is compiled with two \file{pdflatex} passes in
\file{phase2/}. No old cache, git state, README, or existing report is changed.

\section{Limit verdicts and claims that survive}
\begin{longtable}{@{}L{3.1cm}L{7.5cm}L{3.9cm}@{}}
\toprule Limit & Verdict & Check\\\midrule\endhead
Bulk kernel, fixed interior test functions & \C Uniform direct-operator
estimate \eqref{eq:operator}; diagonal layer must be retained. Does not
control the inverse on Bethe densities. & \scr{analytic}\\
WH edge coefficients & \C Corrected formal $A,c,d,e$ survive. Absence of
specified logarithms holds in $u(y)$ through $y^{-2}$. Finite-sea uniform
identification remains \D. & \scr{analytic}, \scr{matched}\\
One-body nested edge $C_T$ & \C $1/(2\pi)$ survives conditionally, with the
independent finite-sea envelope reproduced. Small-momentum-first coefficient
is instead $3/(4\pi)$. & \scr{analytic}, \scr{refit}\\
Fixed nonzero dressed $P$, then small amplitude & \D Not interchangeable by
any result here. Fresh fixed-$\kappa$ scans are diagnostics, not a proof
of the limiting coefficient. & \scr{scan}, \scr{refit}\\
Half-line two-body, $r$ in compact subset of $(0,1)$ & \C Conditional matched
lemma protects $c_0,c_1$ and excludes net $\log^jY/Y$, $j=1,2$, at this order.
Fresh convolution fits support negative $c_1$. & \scr{matched}, \scr{wh},
\scr{refit}\\
Finite sea: $h\to0$ versus $Y\to\infty$ & \D No uniform inverse or differentiated
endpoint bound; the joint overlap requires $Y\gg1$ and $hY\ll1$.
Taking $Y\to\infty$ first at fixed $h$ exits the small-amplitude regime.
& \scr{analytic}, \scr{scan}\\
Equal/vanishing amplitudes & \D No uniform conclusion as $r\to1$ or $r\to0$.
The argument requires both $(1-r^2)Y$ and $r^2Y$ large; constants deteriorate.
& Explicit denominators and domain in matched lemma\\
Physical two-body phase, including head-on & \D The corrected overtaking
convolution does not determine the physical quantum frame, absolute phase,
or a head-on interchange theorem. Earlier audit qualifications remain.
& Scope of \eqref{eq:Q}\\
Phase 2d bound; Phase 2f sign precision & \C The bound remains false and saved
fits do not resolve the subleading sign. Negative half-line signs do not
retroactively make those fits precise. & \scr{refit}\\
\bottomrule
\end{longtable}

\paragraph{What would complete the open part (research judgment, d).}
One needs a uniform bound for the finite-sea resolvent after isolating its
normalisation/zero-frequency mode, with weighted endpoint norms and at least
the derivatives needed for fixed-momentum inversion. The bound must cover an
expanding interval $1\ll Y\ll h^{-1}$ and state how its error grows with $Y$
and with the two amplitude separations. Only then can a joint path
$Y=Y(h)$ be chosen with a proved vanishing error, or compared to the fixed
$\kappa$ semiclassical limit. Physical phase matching additionally requires
the quantum-frame dictionary identified by the controlling audits.

\section{Failures recorded}
\begin{enumerate}
\item \C The historical successive-change assertion $<0.004$ fails at all
three entries reproduced in the numerical section. It is not rerun or
weakened. The new checker asserts that those failures are present in the
saved data. A raw change between two $h$ values is neither a mesh error nor
an extrapolation error bound.
\item \C The Phase 2f sign claim fails under its already-saved $h\log h$
alternative, and correcting the classical curvature alone does not remove
basis dependence. No preferred fit is selected because it agrees with a target.
\item \C The first symbolic check of $e$ failed because SymPy retained
$\psi^{(1)}(-1/2)$ and $\psi^{(2)}(-1/2)$ unevaluated. Explicit recurrence
identities, independently checked at 70 digits, reduce the expression and
pass the unchanged exact-zero assertion. This was a simplification failure,
not evidence that the audited coefficient was wrong. A later exact
classical-limit check also failed on structural equality of factored and
expanded SymPy expressions; simplifying their difference to exact zero
repaired it, with no numerical tolerance introduced.
\item \C The first direct endpoint-moment quadrature at $y=4$ used 12/20
nodes; its 20-node discrepancy $6.22\times10^{-14}$ failed the unchanged
$10^{-14}$ assertion. The grid was refined to 20/32 nodes; the separate
$10^{-10}$ quadrature-spread criterion was also retained.
\item \C The first PDF had a table heading extending 5.51 points beyond its
column. A formatting helper initially used a repository-relative path while
already in \file{phase2/}, causing a file-not-found error. The path was fixed
and the heading shortened before the final two-pass build and visual review.
\end{enumerate}
\C Success markers certify the checks actually asserted, including reproduced
counterexamples. They do not certify the open resolvent assumptions or the
physical continuum identifications.
\end{document}
